% Traduction française COMIMa de l'énoncé original (anglais).
On se donne dans un plan un point $P_0$ et un triangle $A_1A_2A_3$. On pose $A_s = A_{s-3}$ pour tout $s \geq 4$. On construit la suite de points $P_0, P_1, P_2, \ldots$ où $P_{k+1}$ est l'image de $P_k$ par la rotation de centre $A_{k+1}$ et d'angle $120^\circ$ dans le sens des aiguilles d'une montre (sens négatif) $(k = 0, 1, 2, \ldots)$. Montrer que si $P_{1986} = P_0$, alors le triangle $A_1A_2A_3$ est équilatéral (IMO 1986).
